% !TeX root = ../nc-quadric-avatar.tex
% 2026-11-22 / two genuinely distinct genus-one curves in noncubic comparison
\section{Noncubic の点曲線と spectral 曲線の二重 isogeny}
\label{sec:qa-point-spectral-isogeny}

noncubic の中央二重平方環から得る可換点曲線と，
二根対の四次 resolvent から得る spectral 曲線は，
滑らかな場合でさえ一般に同型ではない．
両者の関係を明示することは，avatar の幾何的分類を
内在的 O-marked heart の同値と取り違えないために必要である．

\begin{theorem}[点曲線とspectral曲線の二重 isogeny]
\label{thm:qa-point-spectral-dual-two-isogeny}
$k$ を標数 $0$ の代数閉体，
$f_1,f_2\in k[u,v]_2$ を squarefree かつ互いに
根を共有しない二次形式とする．
\[
 \begin{split}
 E_f&=\{c^2=f_1(u,v),\ d^2=f_2(u,v)\}
                   \ \subset\PP^3_{[u:v:c:d]},\\
 C_q&=\{y^2=f_1(u,v)f_2(u,v)\}
             \ \longrightarrow\PP^1_{[u:v]}
 \end{split}
\]
と置く．$E_f$ は
定理~\cite[\texttt{thm:ore-anticommuting-central}]{QuadricIntrinsic2026}
に現れる中央二重平方環の可換点曲線であり，
$C_q$ は
式~\eqref{eq:qa-rootpair-spectral-curve}の
spectral 曲線である．
次が成り立つ．
\begin{enumerate}
\item $E_f$ と $C_q$ は滑らかな連結種数一曲線である．
 同時符号対合
 $\tau(u,v,c,d)=(u,v,-c,-d)$ は $E_f$ に自由に作用し，
 \[
   r:E_f\longrightarrow C_q,\qquad
   [u:v:c:d]\longmapsto([u:v],y=cd)
   \tag{Point--spectral}\label{eq:qa-point-spectral-cover}
 \]
 はその有限 \'{e}tale 二重商である．
 原点を選べば $r$ は次数二の isogeny となる．
\item $\alpha(c,d)=(-c,d)$ と $\beta(c,d)=(c,-d)$ とする．
 これらの商は滑らかな conic
 \[
    Q_d=\{d^2=f_2(u,v)\},\qquad
    Q_c=\{c^2=f_1(u,v)\}
 \]
 であり，商射を $p_d:E_f\to Q_d$，
 $p_c:E_f\to Q_c$ と書く．
 $Q_d,Q_c\simeq\PP^1$ 上で各々唯一の次数一線束
 $N_d,N_c$ を取り，
 $F_d=p_d^*N_d,\ F_c=p_c^*N_c$ とすると
 \[
   \deg F_d=\deg F_c=2,\quad
   F_d^{\otimes2}\simeq F_c^{\otimes2}
       \simeq\OO_{E_f}(1),\quad
   0\ne P_E:=F_cF_d^{-1}\in\Pic^0(E_f)[2].
   \tag{Point-avatar}\label{eq:qa-point-curve-double-wall}
 \]
 従って $(E_f,F_d,F_c,P_E)$ 自身が
 二重壁上の square avatar を与える．
 また $\tau$ は $P_E$ による非自明二捩れ平行移動である．
\item
 $M_C=(r_*\OO_{E_f})^-$ は
 定理~\ref{thm:qa-noncubic-rootpair-torsionfree-spectral-module}
 の非自明二捩れ線束である．
 その $\Pic^0(C_q)[2]$ の点は，
 $r$ の\emph{双対 isogeny}の核を生成する：
 \[
     0\longrightarrow
       \{\OO_{C_q},M_C\}
       \longrightarrow\Pic^0(C_q)
       \xrightarrow{\ r^*\ }\Pic^0(E_f)
       \longrightarrow0.
       \tag{Dual-2}\label{eq:qa-point-spectral-dual-isogeny}
 \]
 特に原点の選択の下で
 $E_f\simeq C_q/\langle M_C\rangle$ である．
\end{enumerate}
\end{theorem}
\begin{proof}
$E_f\to\PP^1$ は $c,d$ の符号を独立に変える
biquadratic な次数四被覆である．
$f_1,f_2$ の四分岐点が互いに異なるため，
各分岐点では一方の平方根だけが分岐し，
$E_f$ は滑らかで連結である．
Riemann--Hurwitz により
\[
 2g(E_f)-2=4(-2)+4\cdot2=0.
\]
同時符号対合の固定点では$c=d=0$を要するが，
これは$f_1=f_2=0$を意味して仮定に反する．
この対合の不変関数は
$u,v$ と $cd=y$ で，$y^2=f_1f_2$であるから，
商は$C_q$である．
商は非分岐二重被覆で，$C_q$も種数一となる．

$\alpha$ と $\beta$ の不変式から，それぞれの
商が $Q_d,Q_c$ であることが直ちに従う．
squarefree 性によって conic は滑らかであり，
$k$ は代数閉だから $\PP^1$ に同型である．
Conic の超平面線束は $\OO_{\PP^1}(2)$なので，
その平方根となる次数一線束は一意に存在する．
$E_f$ の埋込の線形座標
$(u,v,d)$，$(u,v,c)$ はそれぞれ
$Q_d$，$Q_c$ の超平面切断の引戻しだから
$F_d^2=F_c^2=\OO_{E_f}(1)$となる．
それぞれ次数二で大域生成されており，
相異なる二つの pencil は
補題~\ref{lem:qa-core-two-pencils}により
$E_f$を$\PP^1\times\PP^1$の$(2,2)$因子として
閉埋め込みする．
もし $F_d\simeq F_c$ なら二つの次数二線形系は
基底の射影変換を除いて同じ射を定め，
その deck involution も一致するはずである．
しかし $\alpha\ne\beta$ なので非自明性を得る．

滑らかな種数一曲線の次数二線束 $L$ は，
原点を選んで$x\mapsto \ell_L-x$という
involution を定める．
従って二線束$F_d,F_c$の involution
$\alpha,\beta$の積は，その差
$F_cF_d^{-1}=P_E$による平行移動である
（符号の違いは二捩れゆえ消える）．
これで第二項が従う．

最後に$E_f\to C_q$は \'{e}tale な
$\mu_2$-torsorであり，
$r_*\OO_{E_f}=\OO_{C_q}\oplus M_C$，
$M_C^{\otimes2}\simeq\OO_{C_q}$である．
$H^0(C_q,M_C)=0$は
式~\eqref{eq:qa-rootpair-spectral-module}から従う．
従って $M_C$ は非自明である．
有限 \'{e}tale 同種写像の双対性，
または線束の Galois descent から，
$\ker r^*=\{\OO_{C_q},M_C\}$となる．
次数二の isogeny の双対も次数二なので
$r^*$ は全射となり，最後の完全列を得る．
\end{proof}



\begin{proposition}[非自由MCM二型の行列核は点曲線の二pencilを回収する]
\label{prop:qa-point-mcm-matrix-pencils}
前定理の仮定の下で，$f_1=-\ell_1\ell_2$，
$f_2=-m_1m_2$と因数分解し，
$\ell_1,\ell_2$および$m_1,m_2$を
それぞれ独立な一次形式とする．
中央二重平方環 $B_f$ で
\[
 M_c=\begin{pmatrix}c&\ell_1\\\ell_2&-c\end{pmatrix},
 \qquad
 M_d=\begin{pmatrix}d&m_1\\m_2&-d\end{pmatrix}
 \tag{MCM-matrices}\label{eq:qa-point-mcm-matrices}
\]
と置くと，$M_c^2=M_d^2=0$であり，
それぞれは元のnoncubic normal helixに使う
二つの線形matrix factorizationの表示となる．
$E_f$上では両行列が全点で階数一となり，
その右核の直線から定まる射
\[
 \kappa_c:E_f\longrightarrow\PP^1,\qquad
 \kappa_d:E_f\longrightarrow\PP^1
\]
は，それぞれ前定理の
$p_c:E_f\to Q_c$，$p_d:E_f\to Q_d$
と conic の同型 $Q_c,Q_d\simeq\PP^1$
との合成に一致する（終域の射影変換を除く）．
従って
\[
 \kappa_c^*\OO_{\PP^1}(1)\simeq F_c,\qquad
 \kappa_d^*\OO_{\PP^1}(1)\simeq F_d .
 \tag{MCM-pencils}\label{eq:qa-point-mcm-pencil-classes}
\]
これは\emph{matrix factorizationの点上のkernel}と
幾何的な二pencilの一致であり，
ASのfinite directed End代数全体の同型や
$\pi B_f$を保つpointed心の同型を
まだ結論しない．
\end{proposition}
\begin{proof}
$u,v$ は $B_f$ の中心にあり，
$c^2=f_1=-\ell_1\ell_2$，
$d^2=f_2=-m_1m_2$だから
行列の平方はそれぞれ零である．
これは
\cite[\texttt{thm:ore-anticommuting-central}]{QuadricIntrinsic2026}
の二つの非自由MCM型に用いる
線形行列分解の表示である．
$E_f$の点において$M_c$が零行列となるなら
$c=\ell_1=\ell_2=0$が必要だが，
$\ell_1,\ell_2$は$\PP^1_{[u:v]}$上で
共通零点を持たない．
従って行列の階数は一である．
同じ議論が$M_d$にも適用できる．

$M_c$の成分は$u,v,c$だけに依存するので，
$\kappa_c$は conic $Q_c$ を経由する．
$\ell_1,\ell_2$を $s,t$ とする座標で
$Q_c$の式は $c^2+st=0$である．
その標準パラメータ表示
\[
 [s:t:c]=[a^2:-b^2:ab]
\]
の下で
\[
 M_c=
 \begin{pmatrix}a\\-b\end{pmatrix}
 \begin{pmatrix}b&a\end{pmatrix}
\]
となるから，右核は直線$[a:-b]$である．
これは conicのパラメータ$[a:b]$への
射影線形自己同型であり，
$\kappa_c$は$Q_c\simeq\PP^1$を通る次数二射となる．
よって$\kappa_c^*\OO(1)=F_c$を得る．
$m_1,m_2,d$を用いた同じ計算で
$F_d$の主張も従う．
\end{proof}

\begin{proposition}[MCMの隣接積と二pencilのVeronese積の一致]
\label{prop:qa-point-mcm-composition-veronese}
前命題の仮定の下で，中央二重平方環の
一次部分を
\[
 V=(B_f)_1=\langle u,v,c,d\rangle_k
     \xrightarrow{\ \sim\ }
     H^0(E_f,\OO_{E_f}(1))
 \tag{Degree-one}\label{eq:qa-point-degree-one-restriction}
\]
と同定する．
二つの非自由MCM型 $X_c,X_d$ に関する
noncubic normal helix の隣接合成
\[
 \begin{aligned}
 \mu_c:\;&
 \Hom(\pi X_c,\pi B_f(1))
 \otimes\Hom(\pi B_f,\pi X_c)
 \longrightarrow\Hom(\pi B_f,\pi B_f(1)),\\
 \mu_d:\;&
 \Hom(\pi X_d,\pi B_f(1))
 \otimes\Hom(\pi B_f,\pi X_d)
 \longrightarrow\Hom(\pi B_f,\pi B_f(1))
 \end{aligned}
\]
の像をそれぞれ $U_c,U_d\subset V$ とする．
一方，$E_f$ 上の二つの次数二pencilに対する
通常の対称積
\[
 m_c:\Sym^2H^0(E_f,F_c)\to H^0(E_f,F_c^2),
 \qquad
 m_d:\Sym^2H^0(E_f,F_d)\to H^0(E_f,F_d^2)
\]
の像を前式の$V$へ写す．すると
\[
 \boxed{
 \im\mu_c=\im m_c=U_c=\langle u,v,c\rangle,\qquad
 \im\mu_d=\im m_d=U_d=\langle u,v,d\rangle.}
 \tag{Two-rank3}\label{eq:qa-point-two-rank-three-subspaces}
\]
特に
\[
 \dim U_c=\dim U_d=3,\qquad
 U_c\cap U_d=\langle u,v\rangle,\qquad U_c+U_d=V.
\]
よってnoncubic normal helixの二つのMCM型に伴う
\emph{最小十本矢の隣接合成像}は，点曲線$E_f$の
二つのconic商pencilのVeronese像と一致する．
ただし，この命題は二つの像部分空間と
MCMのkernel写像を同定するもので，
四頂点のfinite directed End代数の
\emph{すべての積}の同型を意味しない．
\end{proposition}
\begin{proof}
$E_f\subset\PP^3$は二quadricの完全交叉で，
一次の斉次関係はない．
よって$V=H^0(E_f,\OO_{E_f}(1))$となる．
定理
~\cite[\texttt{thm:normal-helix} および
 \texttt{cor:normal-cubic-defect}]{QuadricIntrinsic2026}
における非自由MCMの隣接合成は，
線形matrix factorizationの
$2\times2$成分を一次切断とみなしたものだから，
式~\eqref{eq:qa-point-mcm-matrices}より
\[
 \im\mu_c=\langle c,\ell_1,\ell_2\rangle
          =\langle u,v,c\rangle,\qquad
 \im\mu_d=\langle d,m_1,m_2\rangle
          =\langle u,v,d\rangle
\]
である．ここで各一次形式対は独立である．

一方，前命題から
$F_c=p_c^*N_c$，$N_c\simeq\OO_{\PP^1}(1)$であり，
$H^0(E_f,F_c)=H^0(Q_c,N_c)$である．
実際，$p_{c*}\OO_{E_f}$の不変部分は$\OO_{Q_c}$で，
反不変部分の負次数の線束には切断がない．
従って$m_c$は
\[
 \Sym^2H^0(Q_c,N_c)\longrightarrow
 H^0(Q_c,N_c^2)=H^0(Q_c,\OO_{Q_c}(1))
\]
の通常のVeronese積の引戻しである．
この積は三次元切断空間への同型であり，
$Q_c\subset\PP^2_{[u:v:c]}$の
超平面切断の像は
$\langle u,v,c\rangle$である．
同様に$m_d$の像は$\langle u,v,d\rangle$となる．
二つの部分空間の共通部分と和は明らかである．
\end{proof}

\begin{lemma}[noncubic 点曲線への制限と四頂点Endの同型]
\label{lem:qa-pointed-abelianization-four-end}
標数零の代数閉体上、$f_1,f_2$がそれぞれ相異なる二根を持ち、
互いに根を共有しないとする。
$B=B_f$、$R=B/(cd-dc)$、$E=\operatorname{Proj}R=E_f$とする。
$f_1=-\ell_1\ell_2$, $f_2=-m_1m_2$に対する
二つの右$B$-加群$X_c=\operatorname{coker}M_c$,
$X_d=\operatorname{coker}M_d$を前命題のように取る。
$F=F_c$, $G=F_d$と書けば、次の自然な制限は
頂点と全合成を保つ有限次元代数の同型である：
\begin{equation}
\End_{\rm dir}^{\qgr B}
 (\pi B,\pi X_c,\pi B(1),\pi X_d(1))
 \ \xrightarrow{\ \sim\ }\
\End_{\rm dir}^{E}
 (\OO_E,F,F^2,F^2G).
\label{eq:qa-pointed-four-end-abelianization}
\end{equation}
ここで$F^2=G^2=\OO_E(1)$であり、
左辺では$B$の通常シフト$(1)$を記すが、
これ自体を適合Serre平方根とは主張しない。
\end{lemma}
\begin{proof}
$R_c=k[u,v,c]/(c^2-f_1)$とおく。
$R=R_c\oplus dR_c(-1)$は$R_c$上自由である。
滑らかなconic $Q_c=\operatorname{Proj}R_c\simeq\PP^1$
上で$M_c$は$s=\ell_1,t=\ell_2$により
\[
c^2+st=0,\qquad
[s:t:c]=[a^2:-b^2:ab],\qquad
M_c=\binom{a}{-b}(b\ \ a)
\]
となる。従って$\operatorname{coker}M_c$を
$R_c$へ制限した次数付き加群は
$\bigoplus_{n\ge0}H^0(Q_c,\OO_{\PP^1}(2n+1))$である。
$R_c\to R$への自由な拡大と
$p_{c*}\OO_E\simeq\OO_{\PP^1}\oplus\OO_{\PP^1}(-2)$から
\[
X_c\otimes_B R
 \simeq\bigoplus_{n\ge0}H^0(E,F\OO_E(n)),
\qquad
X_d\otimes_B R
 \simeq\bigoplus_{n\ge0}H^0(E,G\OO_E(n)).
\tag{Spinor--sections}\label{eq:qa-pointed-spinor-sections}
\]
両式は単にsheafificationの同型ではなく、
非負次数のsection moduleの同型である。
特に右辺は有限長torsionを持たない。
$B\to R$の核は次数二の$cd-dc$で生成されるので、
$B_0,B_1$ならびに$(X_c)_0,(X_c)_1,
(X_d)_0,(X_d)_1$は、それぞれ$R$への制限で不変である。

Aの適合表示をここで明確にする。
分類論文Aの
\cite[\texttt{thm:ore-anticommuting-central}]{QuadricIntrinsic2026}
で構成するZhang自己同型$\sigma$は、
線形座標変換を伴って$c,d$を交換し、
$B^\sigma$の通常シフトを適合根とする。
右Zhang同値でそのnormal helixの四連続対象を$B$側へ
戻すと、自由加群のshiftの作用を$\sigma$で輸送し、
二つのspinor型も交換するので、
添字を選んで
$(\pi B,\pi X_c,\pi B(1),\pi X_d(1))$
を得る。この輸送は次数付き加群の単なる同一視ではなく、
右作用の$\sigma$-twistを含む。
従って
\cite[\texttt{thm:normal-helix},
\texttt{cor:normal-cubic-defect}]{QuadricIntrinsic2026}
のHom計算と、深さ三によるgraded/tails比較が適用できる。
四頂点の上三角Hom次元は
\[
\begin{pmatrix}
1&2&4&6\\0&1&2&4\\0&0&1&2\\0&0&0&1
\end{pmatrix}.
\tag{Four-End-dims}\label{eq:qa-pointed-four-end-dimensions}
\]
右辺の曲線上の線束も次数$0,2,4,6$であり、
$E$が滑らかな種数一曲線なので同じ次元を持つ。

左辺の各Homは生成元の像で決まり、
$B_0,B_1,(X_d)_0,(X_d)_1$のいずれかに値を持つ。
特に$\Hom(X_c,X_d(1))$の射は
$X_c$の次数零の二生成元の$(X_d)_1$内の像で決まる。
上記の低次数の制限単射性と
\eqref{eq:qa-pointed-spinor-sections}により、
どのHomも右辺のsheaf Homへ単射となる。
両辺の次元が一致するため全て同型である。
$-\otimes_B R$は右完全で、個々の加群射の合成を保つため、
上の六つの同型は同時に全積を保つ。

最初の二段階合成はさらに直接確認できる。
$X_c$の次数零生成元を$e_1,e_2$とし、
$\Hom_B(X_c,B(1))$の基底を
$(c,\ell_1),(-\ell_2,c)$と取ると、
合成の$2\times2$行列は
\[
\begin{pmatrix}
c&\ell_1\\-\ell_2&c
\end{pmatrix}.
\]
従って唯一の二次関係は二つの対角道の差であり、
像は$\langle c,\ell_1,\ell_2\rangle$である。
上のconic表示で行列は
$\binom{a}{b}(b\ \ a)$となり、
$H^0(F)\otimes H^0(F)\to H^0(F^2)$の
反対称核と一致する。この計算は像の一致だけでなく、
二つの二次元Homの\emph{双線形積そのもの}を同定する。
\end{proof}

\begin{theorem}[滑らかなnoncubicの位相を合わせたpointed認識]
\label{thm:qa-point-smooth-noncubic-pointed}
\label{prop:qa-point-mcm-pointed-finite-end-reduction}
前補題の仮定の下で、
square avatar $(E,F,G,P)$を
$F=F_c$, $G=F_d$, $P=GF^{-1}\ne\OO_E$で定める。
$P^2=\OO_E$かつ$G=FP$である。
このsquare型AS$\mathbb Z$-代数を論文Cの指定mutationで
最小十本矢型$A_{\rm av}$に移し、
mutation後のhelixを$(H_i)_{i\in\mathbb Z}$とする。
このとき
\begin{equation}
 \End_{\rm dir}^{\qgr B_f}(E_0,E_1,E_2,E_3)
 \ \simeq\
 \End_{\rm dir}^{\qgr A_{\rm av}}(H_1,H_2,H_3,H_4),
\qquad E_0=\pi B_f,
\label{eq:qa-pointed-shifted-four-end}
\end{equation}
という頂点と全合成を保つ同型が存在する。
論文Aのnormal-helixと本稿のAS10 Serre完成・Noether性・AZ再構成を
適用すると、
\begin{equation}
(\qgr B_f,\pi B_f)\simeq
(\qgr A_{\rm av},\,\pi e_1A_{\rm av})
\tag{Smooth-pointed}\label{eq:qa-pointed-smooth-equivalence}
\end{equation}
というpointed exact同値を得る。
ここで構造対象の像はmutation\emph{後の一番頂点}であり、
ゼロ番頂点ではない。
\end{theorem}
\begin{proof}
square側の曲線上の初期四項は$(P,F,G,FG)$である。
補題~\ref{lem:qa-as10-all-index-mutual-inverse}で扱うevaluation mutationでは
\[
 H_0=K_0,\quad H_1=P,\quad H_2=G,\quad H_3=FG,
\quad
 K_0=\ker\bigl(H^0(FP^{-1})\otimes P\to F\bigr)
      =P^2F^{-1}.
\]
次のsquare項は
系~\ref{cor:qa-core-square-complexes}の式から
$L_4=F^2G^2P,\ L_5=F^3G^2P^2$である。
ゆえに一周期後のevaluation kernelは
\[
 H_4=K_1=L_4^2L_5^{-1}=FG^2.
\]
$P^2=\OO_E$, $G=FP$を代入し、
四項$(H_1,H_2,H_3,H_4)=(P,FP,F^2P,F^3)$を
共通線束$P^{-1}$でtensorすると
\[
(\OO_E,F,F^2,F^3P^{-1})
   =(\OO_E,F,F^2,F^2G).
\]
共通線束のtensorは曲線上の全Homとその積を保つ。
前補題と本稿の全整数mutationのEnd同定を合わせれば
\eqref{eq:qa-pointed-shifted-four-end}を得る。
これは線束上の計算であり、任意のnc heartに
線束tensor自己同値を仮定するものではない。

Aのnormal helixはfullでample、
Cのmutation後のhelixもfull very strongで、
Noether性とAZ amplenessを持つ。
有限Endの同型はtilting Morita同値を与え、
Serre関手と自然に可換する。
従って四項をSerre完成して得られる全整数の
前向きHomと積は一致する。
補題~\ref{lem:qa-as-serre-unique-recovery}と
補題~\ref{lem:qa-as10-all-index-mutual-inverse}および
\cite[\texttt{thm:normal-helix}]{QuadricIntrinsic2026}
の再構成を適用して
\eqref{eq:qa-pointed-smooth-equivalence}を得る。
\end{proof}

\begin{remark}[位相の取り違えと相対境界]
\label{rem:qa-pointed-phase-obstruction}
同じnoncubic normal四項をmutation後の
$(H_0,H_1,H_2,H_3)$に\emph{頂点固定で}
同定することはできない。
前者の$0\to1\to2$合成の核は一次元だが、
後者は最小十本矢型の定義により
$0\to2$に最小関係を持たず、核は零である。
従って一添字shiftは単なる記法でなく必要条件である。
なお定理は四根相異なるsmooth noncubicの
\emph{体上のpointed同値}を述べる。
一根共有・二根対一致における同じ式の相対化、
非被約底上のgroupoid、$C_q$側の通常Picardへの延長、
およびpotential全体の相対忘却は主張しない。
\end{remark}

\begin{corollary}[非零二捩れ付き楕円曲線上の双対二重商]
\label{cor:qa-point-spectral-atkin-lehner}
滑らかな楕円曲線と非零二捩れの組
$(T,L)$ の同型類に対し
\[
 \mathcal W_2(T,L)=
 \bigl(T/\langle L\rangle,\ L^\vee\bigr)
 \tag{Dual-wall}\label{eq:qa-point-spectral-dual-wall}
\]
と置く．ここで $L$ を原点選択によって
位数二の点とみなし，$L^\vee$は
商 isogeny の双対核を指定する非零二捩れである．
この対応は同型類上の involution：
$\mathcal W_2^2=\mathrm{id}$ を定める．
前定理の根対から生じる二つの組は
\[
 \mathcal W_2(C_q,M_C)=(E_f,P_E),\qquad
 \mathcal W_2(E_f,P_E)=(C_q,M_C)
\]
を満たす．
これは楕円曲線と二捩れ\emph{だけ}の対応であり，
$F_d,F_c$と$p^*\OO_{\PP^1}(1)$という
異なる degree-two pencil まで
自動的に同一視するものではない．
\end{corollary}
\begin{proof}
有限 \'{e}tale 二重 isogeny
$T\to T/\langle L\rangle$に対し，
双対 isogeny の核は位数二である．
さらに双対を取れば，楕円曲線の主偏極による
二重双対の自然同型から元の isogeny に戻る．
従って商と双対核の操作は二回で元に戻る．
前定理で$r:E_f\to C_q$の核は$P_E$であり，
$r^\vee:C_q\to E_f$の核は$M_C$と示したので，
二つの表示が得られる．
\end{proof}

\begin{example}[二曲線は一般に同型でない]
\label{ex:qa-point-spectral-j-distinction}
アフィン座標$x=u/v$で
$f_1=x(x-1)$，$f_2=(x-2)(x-3)$とする．
$C_q$の射影模型は
\[
 y^2=x(x-1)(x-2)(x-3).
\]
一方 $E_f$の関数体で
$t=d/(x-3)$と置けば
\[
 x=\frac{3t^2-2}{t^2-1},\qquad
 z=(t^2-1)c,\qquad
 z^2=(3t^2-2)(2t^2-1)=6t^4-7t^2+2.
\]
四次式 $az^4+bz^3+cz^2+dz+e$ の
不変式を $I=12ae-3bd+c^2$，
$J=72ace+9bcd-27ad^2-27b^2e-2c^3$
とすると
$j=1728\cdot4I^3/(4I^3-J^2)$である．
$C_q$の四次式には $(I,J)=(13,-70)$，
$E_f$の上の四次式には $(I,J)=(193,-5362)$を得るので，
\[
 \boxed{j(C_q)=\frac{35152}{9},\qquad
        j(E_f)=\frac{28756228}{3}.}
 \tag{Distinct-j}\label{eq:qa-point-spectral-distinct-j}
\]
従って両曲線は同型でない．
この有理数の計算は
\texttt{computations/verify\_rootpair\_resolvent\_spectral.py}
にも厳密算術の検算を置く．
\end{example}

\begin{remark}[pointed 認識への帰結と退化境界]
\label{rem:qa-point-spectral-pointed-warning}
定理~\ref{thm:qa-noncubic-rootpair-all-geometric-square-avatar-bijection}
で使った square avatar の曲線は \emph{$C_q$} である．
一方，noncubic 中央二重平方環そのものの
可換点曲線は \emph{$E_f$} であり，
第二項は $E_f$ 上にも自然な
二捩れ double-wall avatar があることを示す．
式~\eqref{eq:qa-point-spectral-distinct-j}により，
二つの構成を同一曲線とみなして
avatar由来AS代数と$\qgr B_f$の
pointed同値を結論することはできない．
\emph{どちらの曲線が特定のnormal helixの
AS avatarとして抽出されるか}，
およびそのexceptional foundationの対応は
まだ直接同定していない．これは次の認識問題である．

$f_1,f_2$ が一根を共有すると，
$\tau$には固定点が生じ，
$r$は有限 \'{e}tale 二重isogenyではなくなる．
その際に正しい相対対象は
式~\eqref{eq:qa-rootpair-spectral-module}の
torsion-free $M_C$とその二次積であり，
双対isogenyの短完全列を特異fiberへ
そのまま延長してはならない．
\end{remark}
